\section{Conclusion}

This work demonstrated that the essential motion-control functions of the
SO-101 robot arm can be executed directly on an STM32U545 without a permanent
host-computer connection. The resulting embedded architecture combines direct
Feetech servo communication, calibrated joint limits, encoder feedback,
forward and inverse kinematics, Cartesian feedback control, diagnostics, and
E-Skin supervision in one firmware application. Separation between
CubeMX-generated initialization and handwritten modules kept hardware access,
robot logic, and numerical calculations clearly structured.

The communication and kinematic stack was validated on the physical robot
through checked joint movements and a four-waypoint Cartesian square sequence.
A one-shot IK command alone left a mean model-based TCP error of
\(10.13\,\mathrm{mm}\). The selected Cartesian resolved-rate P controller reduced
this value to \(0.87\,\mathrm{mm}\), with all four reported waypoints below the
\(1\,\mathrm{mm}\) controller criterion. These values show improved consistency
between commanded targets and encoder-based pose estimates, but they do not
establish equivalent physical TCP accuracy or statistical repeatability.

The computational benchmarks showed that optimization of the homogeneous
matrix multiplication was more effective than replacing the standard
trigonometric functions. The selected implementation accelerated complete
forward and inverse kinematics by factors of approximately \(2.07\) and
\(2.02\), respectively, without measured additional error from the matrix
optimization. This confirms that numerical kinematics can be executed within
the resource constraints of the selected microcontroller.

The E-Skin subsystem extended the demonstrator with physically connected
proximity feedback. The ESP32-C6 aggregates measurements and transmits the
maximum value to the STM32, where startup validation, stale-data detection,
motion pause, joint holding, and hysteretic resume behavior are integrated.
However, no systematic measurement of the complete detection-to-hold latency
was performed. The mechanism is therefore a software-based risk-reduction
feature rather than a certified collision-avoidance or emergency-stop system.

Overall, the project establishes a functional standalone platform for further
embedded robot-control development. Its main limitations are position-only
control with four active joints, fixed wrist rotation and gripper state, the
absence of explicit Cartesian path planning, and model-based rather than
independently measured TCP feedback. The logical next step is to combine
repeated trials and external metrology with workspace and singularity analysis,
measured E-Skin response latency, and explicit Cartesian trajectory generation.
Full-pose control and any safety-related claim would require additional sensing,
hardware safeguards, and a dedicated validation process.
